%%%% packages %%%%

% input and output encodings
\usepackage[utf8]{inputenc} % input font encoding that supports special characters
\usepackage[T1]{fontenc} % output font encoding that supports special characters

\setlength{\parindent}{0pt}

% language settings
% German-specific commands
\usepackage[ngerman]{babel}
% English-specific commands
%\usepackage[english]{babel}

\usepackage[bottom]{footmisc} % left aligned footnotes / multiple footnotes

% hyper-links
\usepackage[colorlinks = true, % colored hyper-links
			pagebackref = true, % allows to add links in the bibliography
			breaklinks = true, % allows for multi-lined hyper-links
            urlcolor = blue, % color of external hyper-links
            linkcolor = red, % color of internal hyper-links
]{hyperref}

\hypersetup{linkcolor=black}

% graphics
\usepackage{graphicx} % to include graphics
\usepackage{sidecap} %  defines environments called SCfigure and SCtable (analogous to figure and table) to typeset captions sideways.
\usepackage{caption} %  provides many ways to customise the captions in floating environments

% drawing
\usepackage{tikz} % to draw figures with tikz

\usepackage{verbatim} % \begin{comment}...\end{comment}

% colors
\usepackage{xcolor} % provides several kinds of color tints shades, tones etc.

% enumerations
\usepackage{enumerate} % adds an optional argument to the enumerate environment which determines the style in which the counter is printed

% mathematics
\usepackage{amsmath} % miscellaneous enhancements for structure of documents containing mathematical formulas
\usepackage{amssymb} % extended collection of mathematical symbols
\usepackage{amscd} % mathematical commutative arrow diagrams
\usepackage{amsthm} % environments for theorems and lemmas etc.
\usepackage{bm} % bold math

% spacing
\usepackage{lscape} % landscape pages in a portrait document
\usepackage{setspace} % spaces between lines

% listings
\usepackage{minted} % to insert code

% quotation
\usepackage{csquotes} % quoting environment

\usepackage{enumitem} % provides user control over the layout of the three basic list environments: enumerate, itemize and description

\usepackage[framemethod=TikZ]{mdframed} % develops the facilities of framed in providing breakable framed and coloured boxes

\setcounter{chapter}{0} % to start with chapter 0 not chapter 1

\usepackage{tocbibind} % include list of figures, tables and bibliography in the table of contents (toc)

% to modify the position of the page numbering
\usepackage{fancyhdr} % to modify headers
\fancyhf{} % sets both header and footer to nothing
%\renewcommand{\headrulewidth}{0pt}
%\pagestyle{fancy}

\usepackage{lmodern} % The Latin Modern fonts are enhanced versions of the Computer Modern fonts.

\usepackage[a4paper,bindingoffset=0.5cm,%
            left=2.5cm,right=2.5cm,top=2.5cm,bottom=2.5cm,%
            footskip=1.5cm]{geometry} % wider text


\usepackage[lastexercise,answerdelayed]{exercise}
\renewcommand{\ExerciseName}{Exercice}
\renewcommand{\AnswerName}{Exercice}
\usepackage{totcount}
\regtotcounter{Exercise} % register the counter for getting the total
\usepackage{chngcntr}
\counterwithin{Exercise}{chapter}
\counterwithin{Answer}{chapter}
\usepackage{xassoccnt}
\NewTotalDocumentCounter{totalex}
\usepackage{titleref}
\usepackage{etoolbox}
\renewcommand{\ExerciseHeader}{ \stepcounter{totalex}\ifnumcomp{\value{Exercise}}{=}{1}{\ifnumcomp{\thechapter}{=}{1}{}{\vspace{10pt}}\noindent\par\vspace{10pt}}{}\noindent\normalsize\bfseries\ExerciseName\ \ExerciseHeaderNB\ \ExerciseHeaderDifficulty\par\medskip}
\renewcommand{\AnswerHeader}{\ifnumcomp{\value{Exercise}}{=}{1}{\ifnumcomp{\thechapter}{=}{1}{}{\vspace{10pt}}\noindent\par\vspace{10pt}}{}\noindent\normalsize\bfseries\AnswerName\ \ExerciseHeaderNB\ \ExerciseHeaderDifficulty\par\medskip}


\usepackage[framemethod=TikZ]{mdframed}
%%%%%%%%%%%%%%%%%%%%%%%%%%%%%%
%Theorem
\newcounter{theo}[chapter] \setcounter{theo}{0}
\renewcommand{\thetheo}{\arabic{chapter}.\arabic{theo}}
\newenvironment{theo}[2][]{%
\refstepcounter{theo}%
\ifstrempty{#1}%
{\mdfsetup{%
frametitle={%
\tikz[baseline=(current bounding box.east),outer sep=0pt]
\node[anchor=east,rectangle,fill=blue!20]
{\strut Satz~\thetheo};}}
}%
{\mdfsetup{%
frametitle={%
\tikz[baseline=(current bounding box.east),outer sep=0pt]
\node[anchor=east,rectangle,fill=blue!20]
{\strut Satz~\thetheo:~#1};}}%
}%
\mdfsetup{innertopmargin=10pt,linecolor=blue!20,%
linewidth=2pt,topline=true,%
frametitleaboveskip=\dimexpr-\ht\strutbox\relax
}
\begin{mdframed}[]\relax%
\label{#2}}{\end{mdframed}}
%%%%%%%%%%%%%%%%%%%%%%%%%%%%%%
%Lemma
\newcounter{lem}[chapter]\setcounter{lem}{0}
\renewcommand{\thelem}{\arabic{chapter}.\arabic{lem}}
\newenvironment{lem}[2][]{%
\refstepcounter{lem}%
\ifstrempty{#1}%
{\mdfsetup{%
frametitle={%
\tikz[baseline=(current bounding box.east),outer sep=0pt]
\node[anchor=east,rectangle,fill=green!20]
{\strut Lemma~\thelem};}}
}%
{\mdfsetup{%
frametitle={%
\tikz[baseline=(current bounding box.east),outer sep=0pt]
\node[anchor=east,rectangle,fill=green!20]
{\strut Lemma~\thetheo:~#1};}}%
}%
\mdfsetup{innertopmargin=10pt,linecolor=green!20,%
linewidth=2pt,topline=true,%
frametitleaboveskip=\dimexpr-\ht\strutbox\relax
}
\begin{mdframed}[]\relax%
\label{#2}}{\end{mdframed}}

%%%%%%%%%%%%%%%%%%%%%%%%%%%%%%
%Definition
\newcounter{deff}[chapter]\setcounter{deff}{0}
\renewcommand{\thedeff}{\arabic{chapter}.\arabic{deff}}
\newenvironment{deff}[2][]{%
\refstepcounter{deff}%
\ifstrempty{#1}%
{\mdfsetup{%
frametitle={%
\tikz[baseline=(current bounding box.east),outer sep=0pt]
\node[anchor=east,rectangle,fill=yellow!40]
{\strut Definition~\thedeff};}}
}%
{\mdfsetup{%
frametitle={%
\tikz[baseline=(current bounding box.east),outer sep=0pt]
\node[anchor=east,rectangle,fill=yellow!40]
{\strut Definition~\thedeff:~#1};}}%
}%
\mdfsetup{innertopmargin=10pt,linecolor=yellow!40,%
linewidth=2pt,topline=true,%
frametitleaboveskip=\dimexpr-\ht\strutbox\relax
}
\begin{mdframed}[]\relax%
\label{#2}}{\end{mdframed}}

%%%%%%%%%%%%%%%%%%%%%%%%%%%%%%
%Proof
\newcounter{prf}[chapter]\setcounter{prf}{0}
\renewcommand{\theprf}{\arabic{chapter}.\arabic{prf}}
\newenvironment{prf}[2][]{%
\refstepcounter{prf}%
\ifstrempty{#1}%
{\mdfsetup{%
frametitle={%
\tikz[baseline=(current bounding box.east),outer sep=0pt]
\node[anchor=east,rectangle,fill=red!20]
{\strut Beweis~\theprf};}}
}%
{\mdfsetup{%
frametitle={%
\tikz[baseline=(current bounding box.east),outer sep=0pt]
\node[anchor=east,rectangle,fill=red!20]
{\strut Beweis~\theprf:~#1};}}%
}%
\mdfsetup{innertopmargin=10pt,linecolor=red!20,%
linewidth=2pt,topline=true,%
frametitleaboveskip=\dimexpr-\ht\strutbox\relax
}
\begin{mdframed}[]\relax%
\label{#2}}{\qed\end{mdframed}}

%%%%%%%%%%%%%%%%%%%%%%%%%%%%%%
%Example
\newcounter{example}[chapter]\setcounter{example}{0}
\renewcommand{\theexample}{\arabic{chapter}.\arabic{example}}
\newenvironment{example}[2][]{%
\refstepcounter{example}%
\ifstrempty{#1}%
{\mdfsetup{%
frametitle={%
\tikz[baseline=(current bounding box.east),outer sep=0pt]
\node[anchor=east,rectangle,fill=gray!20]
{\strut Beispiel~\theexample};}}
}%
{\mdfsetup{%
frametitle={%
\tikz[baseline=(current bounding box.east),outer sep=0pt]
\node[anchor=east,rectangle,fill=gray!20]
{\strut Beispiel~\theexample:~#1};}}%
}%
\mdfsetup{innertopmargin=10pt,linecolor=gray!20,%
linewidth=2pt,topline=true,%
frametitleaboveskip=\dimexpr-\ht\strutbox\relax
}
\begin{mdframed}[]\relax%
\label{#2}}{\end{mdframed}}


% empty set
\let\oldemptyset\emptyset
\let\emptyset\varnothing

\usepackage{mathtools} % for ceil and floor




% tikz graphs
\tikzstyle{vertex}=[circle,fill=black!25,minimum size=20pt,inner sep=0pt]
\tikzstyle{selected vertex} = [vertex, fill=red!24]
\tikzstyle{edge} = [draw,thick,-]
\tikzstyle{weight} = [font=\small]
\tikzstyle{selected edge} = [draw,line width=5pt,-,red!50]
\tikzstyle{ignored edge} = [draw,line width=5pt,-,black!20]

\usepackage{pgf}
\usetikzlibrary{automata, positioning, arrows}

\usetikzlibrary{chains,fit,shapes}

% define symbol :=
\mathchardef\ordinarycolon\mathcode`\:
\mathcode`\:=\string"8000
\begingroup \catcode`\:=\active
  \gdef:{\mathrel{\mathop\ordinarycolon}}
\endgroup


% more text colors
%\usepackage[dvipsnames]{xcolor}


\newlength{\tmtapesize}
\setlength{\tmtapesize}{9mm}
\tikzset{tmtape/.style={
        draw,
        minimum size=\tmtapesize,
        inner sep=0,
        text height=1.8ex,
        text depth=0.25ex,
        on chain
    }
}
\tikzset{tmtapehend/.style={
        minimum size=\tmtapesize,
        line width=0pt,
        on chain,
        append after command={
            (\tikzlastnode.north west) edge (\tikzlastnode.north east)
            (\tikzlastnode.south west) edge (\tikzlastnode.south east)
        }
    }
}
\tikzset{tmtapevend/.style={
        minimum size=\tmtapesize,
        line width=0pt,
        on chain,
        append after command={
            (\tikzlastnode.north west) edge (\tikzlastnode.south west)
            (\tikzlastnode.north east) edge (\tikzlastnode.south east)
        }
    }
}
\tikzset{tmtapecursor/.style={
        minimum size=\dimexpr\tmtapesize+2mm,
        draw=black,
        rounded corners=1mm,
        fill=blue,
        fill opacity=0.1
    }
}

